\chapter{Persamaan Diferensial Linear}\label{ch:b1:diffeq}

Pertama kali ditemui pada jilid sekolah menengah, persamaan diferensial di sini
diperlakukan dengan bukti lengkap dan dalam keumuman yang lebih besar: persamaan
linear orde satu berkoefisien variabel (yang diselesaikan tuntas dengan
metode \omterm{thm:b1:diffeq:voc}{variasi konstanta}) dan \omterm{def:b1:diffeq:linear2}{persamaan linear orde dua
berkoefisien konstan}, yaitu model bagi osilasi. Kedua kasusnya memperlihatkan
struktur yang sama: \emph{penyelesaian umum $=$ satu penyelesaian khusus
$+$ penyelesaian umum persamaan homogennya}.

\section{Persamaan linear orde satu}

\begin{definition}\label{def:b1:diffeq:linear1}
Misalkan $I$ sebuah selang dan $a, b \colon I \to \R$ (atau $\C$)
kontinu. Persamaan
\[
(E)\colon\quad y' + a(x)\,y = b(x),
\]
dengan fungsi $y \colon I \to \R$ (atau $\C$) yang dapat diturunkan sebagai variabelnya,
disebut \emph{persamaan diferensial linear orde
satu}\index{persamaan diferensial!linear orde satu}. Persamaan
$(H)\colon y' + a(x) y = 0$ adalah persamaan \emph{homogen}
yang bersesuaian dengannya.
\end{definition}

\begin{theorem}[Menyelesaikan persamaan homogennya]\label{thm:b1:diffeq:homogeneous1}
Misalkan $A$ sebuah antiturunan $a$ pada $I$ (yang memang ada:
\cref{ch:b1:integration}). Penyelesaian $(H)$ pada $I$ tepat
berupa fungsi
\[
y(x) = \lambda\, \eu^{-A(x)}, \qquad \lambda \in \R \text{ (atau } \C).
\]
\end{theorem}

\begin{proof}
Fungsi itu memang penyelesaiannya: $y' = -\lambda A' \eu^{-A} = -a y$.
Sebaliknya, misalkan $y$ memenuhi $(H)$ lalu tulis $z(x) = y(x)\, \eu^{A(x)}$.
Maka
\[
z' = y' \eu^{A} + y\, a\, \eu^{A} = (y' + ay)\, \eu^{A} = 0 ,
\]
sehingga $z$ konstan pada selang $I$, katakanlah $z = \lambda$: jadi $y =
\lambda \eu^{-A}$. (Perhatikan logikanya: tak ada penyelesaian yang hilang, karena
\emph{setiap} penyelesaian sudah dituliskan dalam bentuk yang dinyatakan tadi.)
\end{proof}

\begin{example}[Persamaan homogen berkoefisien variabel]\label{ex:b1:diffeq:varcoeff}
Selesaikan $y' + (\cos x)\,y = 0$ pada $\R$. Sebuah antiturunan $a(x) =
\cos x$ adalah $A(x) = \sin x$, jadi penyelesaiannya adalah
\[
y(x) = \lambda\,\eu^{-\sin x}, \qquad \lambda \in \R .
\]
Dua pembacaan. Setiap penyelesaiannya periodik (dengan periode $2\pi$) dan tak pernah
lenyap kecuali bila $\lambda = 0$ --- tanda $\lambda$ adalah
tanda $y$ untuk selamanya, karena eksponensial tak dapat memotong nol. Dan
penyelesaian yang melewati $y(0) = y_0$ adalah $y_0\eu^{-\sin x}$: jadi tepat
satu kurva keluarga itu melewati setiap titik awal, yaitu gambaran
berdimensi satu dari \cref{thm:b1:diffeq:voc}~(2).
\end{example}

\begin{theorem}[Variasi konstanta; masalah Cauchy]\label{thm:b1:diffeq:voc}
Dengan notasi di atas:
\begin{enumerate}
  \item Penyelesaian $(E)$ pada $I$ tepat berupa
        \[
        y(x) = \Bigl(\lambda + \int_{x_0}^{x} b(t)\, \eu^{A(t)} \dd t
        \Bigr)\, \eu^{-A(x)},
        \qquad \lambda \in \R,
        \]
        dengan $x_0 \in I$ yang tetap. Setara dengan itu: penyelesaian umum
        $(H)$ ditambah satu penyelesaian khusus $(E)$.
  \item Untuk setiap $x_0 \in I$ dan $y_0$, \emph{masalah
        Cauchy}\index{masalah Cauchy} ``$(E)$ beserta $y(x_0) = y_0$''
        mempunyai tepat satu penyelesaian pada $I$.
\end{enumerate}
\end{theorem}

\begin{proof}
(1) Dengan mengikuti metode yang bernama \emph{variasi
konstanta}\index{variasi konstanta}, cari penyelesaian yang
berbentuk $y = \mu(x)\, \eu^{-A(x)}$ dengan $\mu$ dapat diturunkan --- tak ada
keumuman yang hilang, karena setiap fungsi pada $I$ dapat ditulis demikian
($\mu = y\,\eu^{A}$). Dengan mensubstitusikannya,
\[
y' + ay = \mu' \eu^{-A} - \mu a \eu^{-A} + a\mu\eu^{-A}
= \mu'\, \eu^{-A},
\]
sehingga $y$ memenuhi $(E)$ jika dan hanya jika $\mu'(x) = b(x)\,\eu^{A(x)}$, yakni
jika dan hanya jika $\mu(x) = \lambda + \int_{x_0}^x b(t)\eu^{A(t)}\dd t$ untuk
suatu konstanta $\lambda$ (karena dua antiturunan dari fungsi kontinu yang sama
pada sebuah selang berselisih sebuah konstanta).

(2) Pada rumus itu, $y(x_0) = \lambda\,\eu^{-A(x_0)}$: syarat
$y(x_0) = y_0$ menentukan $\lambda = y_0 \eu^{A(x_0)}$ secara tunggal.
\end{proof}

\begin{example}\label{ex:b1:diffeq:order1}
Selesaikan $y' + \dfrac{y}{x} = x^2$ pada $I = \intoo{0}{+\infty}$. Di sini
$a(x) = \frac 1x$, $A(x) = \ln x$, $\eu^{-A(x)} = \frac 1x$.
Penyelesaian homogennya: $\frac{\lambda}{x}$. \omterm{thm:b1:diffeq:voc}{Variasi konstanta}:
$\mu'(x) = x^2 \cdot x = x^3$, jadi $\mu = \frac{x^4}{4} + \lambda$, dan
\[
y(x) = \frac{x^3}{4} + \frac{\lambda}{x}, \qquad \lambda \in \R .
\]
Dengan syarat awal $y(1) = 0$: $\lambda = -\frac14$.
\emph{Periksa:} $y' + \frac yx = \frac{3x^2}{4} - \frac{\lambda}{x^2} +
\frac{x^2}{4} + \frac{\lambda}{x^2} = x^2$.
\end{example}

\begin{example}[Menerka mengalahkan mengintegralkan]\label{ex:b1:diffeq:guess}
Selesaikan $y' + 2x\,y = x$ pada $\R$. \omterm{thm:b1:diffeq:voc}{Variasi konstanta} memang berhasil ($A =
x^2$, $\mu' = x\,\eu^{x^2}$, $\mu = \frac12\eu^{x^2} + \lambda$),
tetapi mengamati bahwa \emph{konstanta} $y_p = \frac12$ memenuhi
persamaan itu ($0 + 2x\cdot\frac12 = x$) jauh lebih cepat. Dengan
penyelesaian homogennya $\lambda\,\eu^{-x^2}$:
\[
y(x) = \frac12 + \lambda\,\eu^{-x^2}, \qquad \lambda \in \R .
\]
Setiap penyelesaiannya konvergen ke $\frac12$ dengan sangat cepat ketika $x \to
\pm\infty$: penyelesaian khusus yang konstan itu adalah sebuah
\emph{kesetimbangan} yang dihampiri semua penyelesaian. Inti gagasannya: sebelum
meluncurkan metode umumnya, luangkan sepuluh detik untuk mencari
penyelesaian khusus yang kasatmata (konstanta, monomial, kelipatan
ruas kanan); teorema strukturnya lalu merampungkan pekerjaannya.
\end{example}

\begin{remark}[Selangnya penting]
Teorema itu hidup pada sebuah \emph{selang} yang di situ $a$ dan $b$
kontinu. Untuk $y' + \frac yx = 0$ pada $\R^*$, penyelesaiannya adalah
$\frac{\lambda}{x}$ pada $\intoo{0}{+\infty}$ dan
$\frac{\mu}{x}$ pada $\intoo{-\infty}{0}$ dengan konstanta yang
\emph{saling bebas}: tak ada alasan bagi satu rumus tunggal untuk merekat melewati
kesingularan di $0$.
\end{remark}

\begin{example}[Ruas kanan yang kompleks, dua jawaban real]\label{ex:b1:diffeq:complexrhs}
Selesaikan $y' - y = \cos x$ dan $y' - y = \sin x$ sekaligus. Bekerjalah
di $\C$ dengan ruas kanan $\eu^{\iu x}$: mencoba $y_p =
c\,\eu^{\iu x}$ memberikan $c(\iu - 1)\eu^{\iu x} = \eu^{\iu x}$, jadi
\[
c = \frac1{\iu - 1} = \frac{-1 - \iu}2,
\qquad
y_p = -\frac{(1 + \iu)(\cos x + \iu\sin x)}2
= \frac{\sin x - \cos x}2 + \iu\,\frac{-\sin x - \cos x}2 .
\]
Karena persamaannya berkoefisien real, bagian real dan bagian
imajinernya terpisah: $\frac{\sin x - \cos x}2$ memenuhi $y' - y = \cos x$,
dan $-\frac{\sin x + \cos x}2$ memenuhi $y' - y = \sin x$ (periksa
yang pertama: turunannya $\frac{\cos x + \sin x}2$, dikurangi
fungsinya, memberikan $\cos x$). Satu baris kompleks menggantikan dua kali
\omterm{thm:b1:diffeq:voc}{variasi konstanta} --- penghematan yang sama seperti yang disistematiskan
\cref{met:b1:diffeq:particular} untuk orde dua,
dan dividen yang berulang dari \cref{ch:b1:complex}.
\end{example}

\section{Persamaan linear orde dua berkoefisien konstan}

\begin{definition}\label{def:b1:diffeq:linear2}
Misalkan $a, b \in \R$ dan $f \colon I \to \R$ kontinu. Persamaan
\[
(E)\colon\quad y'' + a\,y' + b\,y = f(x)
\]
disebut \emph{persamaan linear orde dua berkoefisien
konstan}\index{persamaan diferensial!linear orde dua};
$(H)\colon y'' + ay' + by = 0$ adalah persamaan homogennya, dan
$\chi(r) = r^2 + ar + b$ adalah \emph{polinomial
karakteristik}\index{polinomial karakteristik}nya.
\end{definition}

\begin{theorem}[Penyelesaian homogennya]\label{thm:b1:diffeq:homogeneous2}
Misalkan $\Delta = a^2 - 4b$ diskriminan $\chi$. Penyelesaian real
$(H)$ pada $\R$ adalah:
\begin{enumerate}
  \item jika $\Delta > 0$, dengan $r_1 \neq r_2$ kedua akar realnya:
        $\;y = \lambda\, \eu^{r_1 x} + \mu\, \eu^{r_2 x}$;
  \item jika $\Delta = 0$, dengan $r_0$ akar gandanya:
        $\;y = (\lambda + \mu x)\, \eu^{r_0 x}$;
  \item jika $\Delta < 0$, dengan akar $\alpha \pm \iu\omega$
        ($\omega > 0$):
        $\;y = \eu^{\alpha x} (\lambda \cos\omega x + \mu
        \sin\omega x)$;
\end{enumerate}
dan pada masing-masing kasus $(\lambda, \mu)$ menjelajahi $\R^2$.
\end{theorem}

\begin{proof}
Amati lebih dulu bahwa untuk $r \in \C$, $x \mapsto \eu^{rx}$ memenuhi $(H)$
jika dan hanya jika $\chi(r) = 0$ (substitusikan: $(r^2 + ar + b)\eu^{rx} =
0$). Itulah sebabnya eksponensial menjadi terkaan pertama yang wajar:
pendiferensialan bekerja pada $\eu^{rx}$ sebagai perkalian dengan bilangan
$r$, sehingga persamaan diferensialnya menjadi persamaan numerik
$\chi(r) = 0$ --- seluruh masalah analitisnya termampatkan menjadi
pencarian akar satu persamaan kuadrat.

Langkah kuncinya adalah pergantian variabel yang menurunkan ordenya. Misalkan $r$
sebuah akar $\chi$ (yang mungkin kompleks) lalu tulis $y = z\, \eu^{rx}$,
yang tidak menghilangkan keumuman. Maka
\[
y'' + ay' + by
= \bigl(z'' + (2r + a) z' + \chi(r) z\bigr)\eu^{rx}
= \bigl(z'' + (2r + a) z'\bigr)\eu^{rx},
\]
sehingga $(H)$ menjadi persamaan \emph{orde satu} $u' + (2r + a) u = 0$
bagi $u = z'$.

\emph{Kasus $\Delta \neq 0$:} pilih $r = r_1$; maka $2r_1 + a = r_1 -
r_2$ (karena $r_1 + r_2 = -a$). Menurut
\cref{thm:b1:diffeq:homogeneous1}, $z' = c\,\eu^{(r_2 - r_1)x}$ untuk
suatu konstanta $c$; dengan mengintegralkannya pada $\R$, $z = \mu\, \eu^{(r_2 - r_1)x}
+ \lambda$ dengan $\mu = \frac{c}{r_2 - r_1}$, sehingga $y = z\,
\eu^{r_1 x} = \lambda\,\eu^{r_1x} + \mu\,\eu^{r_2x}$.
Ketika $\Delta < 0$, akarnya adalah $\alpha \pm \iu\omega$ dan
penyelesaian kompleksnya adalah $y = c_1\eu^{(\alpha+\iu\omega)x} +
c_2\eu^{(\alpha-\iu\omega)x}$ dengan $c_1, c_2 \in \C$. Yang mana di antaranya
yang bernilai real? Karena $\conj{\eu^{(\alpha+\iu\omega)x}} =
\eu^{(\alpha-\iu\omega)x}$, \omterm{def:b1:complex:field}{konjugat} $y$ adalah $\conj{c_2}\,
\eu^{(\alpha+\iu\omega)x} + \conj{c_1}\,\eu^{(\alpha-\iu\omega)x}$,
dan $y = \conj y$ untuk setiap $x$ memaksa $c_2 = \conj{c_1}$ (karena kedua
eksponensial itu bebas linear: hitung nilainya di dua titik, atau
bandingkan di $x=0$ setelah membaginya dengan $\eu^{\alpha x}$). Dengan menulis $c_1
= \frac{\lambda - \iu\mu}2$ dengan $\lambda, \mu$ real:
\[
y = 2\,\Re\Bigl(c_1\,\eu^{\alpha x}(\cos\omega x +
\iu\sin\omega x)\Bigr)
= \eu^{\alpha x}\bigl(\lambda\cos\omega x + \mu\sin\omega
x\bigr),
\]
dan sebaliknya setiap fungsi semacam itu merupakan penyelesaian (yaitu bagian real
sebuah penyelesaian kompleks bagi persamaan yang real): jadi ruang penyelesaian realnya
persis seperti yang dinyatakan tadi.

\emph{Kasus $\Delta = 0$:} $r = r_0$, $2r_0 + a = 0$, jadi $z'' = 0$:
sehingga $z = \lambda + \mu x$ dan $y = (\lambda + \mu x)\eu^{r_0 x}$.
\end{proof}

\begin{example}[Sebuah masalah Cauchy, dari awal sampai akhir]\label{ex:b1:diffeq:cauchy2}
Selesaikan $y'' - 3y' + 2y = 0$ dengan $y(0) = 0$, $y'(0) = 1$.
\omterm{def:b1:diffeq:linear2}{Polinomial karakteristiknya} $r^2 - 3r + 2 = (r - 1)(r - 2)$ mempunyai
akar real $1$ dan $2$: jadi penyelesaian umumnya $y = \lambda\eu^{x} +
\mu\eu^{2x}$. Kedua syaratnya memberikan sistem linear
\[
\lambda + \mu = 0, \qquad \lambda + 2\mu = 1 ,
\]
jadi $\mu = 1$, $\lambda = -1$:
\[
y(x) = \eu^{2x} - \eu^{x} .
\]
Periksa: $y(0) = 0$; $y' = 2\eu^{2x} - \eu^x$ memberikan $y'(0) = 1$; dan
$y'' - 3y' + 2y = (4 - 6 + 2)\eu^{2x} + (-1 + 3 - 2)\eu^x = 0$.
Perhatikan bentuk jawabannya: di dekat $-\infty$ ragam yang lambat
$-\eu^x$ mendominasi; di dekat $+\infty$ ragam yang cepat $\eu^{2x}$ yang
mendominasi. Membaca penyelesaian sebagai superposisi ragam dengan laju peluruhan
atau pertumbuhan yang berbeda adalah kebiasaan yang menguntungkan --- begitulah
pemilahan peralihan dan keadaan tunak pada soal akhir pekan ditata.
\end{example}

\begin{omfigure}
\begin{tikzpicture}[xscale=1.35, yscale=1.55]
  \draw[->, thin] (0,0) -- (6.4,0) node[below] {$x$};
  \draw[->, thin] (0,-0.85) -- (0,1.25) node[left] {$y$};
  \draw[omDef, thick, domain=0:6.2, samples=400, smooth]
    plot (\x, {exp(-0.35*\x)*cos(4*\x r)});
  \draw[omThm, thick, domain=0:6.2, samples=120, smooth]
    plot (\x, {(1+1.8*\x)*exp(-1.1*\x)});
  \draw[omProp, thick, domain=0:6.2, samples=120, smooth]
    plot (\x, {1.15*exp(-0.45*\x) - 0.15*exp(-2.5*\x)});
  \node[omDef, font=\small] at (2.1,-0.62)
    {$\Delta < 0$: berosilasi};
  \node[omThm, font=\small] at (0.95,1.1) {$\Delta = 0$};
  \node[omProp, font=\small] at (4.9,0.38)
    {$\Delta > 0$: tanpa osilasi};
\end{tikzpicture}

{\small Ketiga regim $y'' + ay' + by = 0$ yang penyelesaiannya
meluruh: osilasi teredam (akar kompleks), kembali kritis
(akar ganda), peluruhan teredam lebih (dua akar real). Regim mana yang
terjadi terbaca dari tanda $\Delta = a^2 - 4b$ saja ---
sebelum apa pun diselesaikan.}
\end{omfigure}

\begin{theorem}[Struktur dan masalah Cauchy]\label{thm:b1:diffeq:structure2}
\begin{enumerate}
  \item Jika $y_p$ satu penyelesaian khusus $(E)$, maka penyelesaian
        $(E)$ tepat berupa $y_p + y_h$, dengan $y_h$ menjelajahi
        penyelesaian $(H)$.
  \item (Superposisi) Jika $y_1$ memenuhi $y'' + ay' + by = f_1$ dan
        $y_2$ memenuhi $y'' + ay' + by = f_2$, maka $y_1 + y_2$ memenuhi
        persamaan dengan ruas kanan $f_1 + f_2$.
  \item Untuk setiap $x_0 \in I$ dan $(y_0, y_0')$, \omterm{thm:b1:diffeq:voc}{masalah Cauchy}
        ``$(E)$, $y(x_0) = y_0$, $y'(x_0) = y_0'$'' mempunyai tepat satu
        penyelesaian pada $I$. \emph{(Keberadaannya diberikan oleh sebuah
        penyelesaian khusus; ketunggalannya penuh.)}
\end{enumerate}
\end{theorem}

\begin{proof}
(1) Di sini $y$ memenuhi $(E)$ jika dan hanya jika $y - y_p$ memenuhi $(H)$, menurut kelinearan $y
\mapsto y'' + ay' + by$. Butir (2) adalah kelinearan yang sama.

(3) Menurut (1) cukup dibuktikan bahwa konstanta $(\lambda,
\mu)$ selalu dapat disetel, secara tunggal, ke sebarang data $(y_0, y_0')$.
Dengan menggeser variabelnya, andaikan $x_0 = 0$. Pada kasus (1)
\cref{thm:b1:diffeq:homogeneous2}, $y = \lambda\eu^{r_1x} +
\mu\eu^{r_2x}$ gives
\[
y(0) = \lambda + \mu, \qquad y'(0) = r_1\lambda + r_2\mu :
\]
sebuah sistem linear dalam $(\lambda, \mu)$ yang determinannya $r_2 -
r_1 \neq 0$; dengan menyelesaikannya secara eksplisit, $\mu = \frac{y_0' -
r_1y_0}{r_2 - r_1}$ dan $\lambda = y_0 - \mu$: jadi tepat satu
penyelesaian. Pada kasus (2), $y(0) = \lambda$ dan $y'(0) = r_0\lambda +
\mu$: sistemnya segitiga dengan determinan $1$, yang diselesaikan oleh
$\lambda = y_0$, $\mu = y_0' - r_0y_0$. Pada kasus (3), $y(0) =
\lambda$ dan $y'(0) = \alpha\lambda + \omega\mu$: determinannya
$\omega \neq 0$, yang diselesaikan oleh $\lambda = y_0$, $\mu = \frac{y_0' -
\alpha y_0}\omega$. Pada masing-masing kasus \omterm{def:b1:logic:map}{pemetaan} $(\lambda, \mu) \mapsto
(y(x_0), y'(x_0))$ adalah bijeksi linear --- dan bahasa
\cref{ch:b1:linmaps} akan memampatkan pemeriksaan kasus ini menjadi satu
kalimat.
\end{proof}

\begin{method}[Penyelesaian khusus untuk $f(x) = P(x)\,\eu^{\gamma x}$]\label{met:b1:diffeq:particular}
Ketika ruas kanannya $P(x)\,\eu^{\gamma x}$ dengan $P$ sebuah
polinomial dan $\gamma \in \R$ (yang mencakup polinomial, eksponensial,
dan lewat $\gamma$ kompleks atau superposisi, juga $\cos$ dan $\sin$): carilah
penyelesaian khusus yang berbentuk
\[
y_p(x) = x^{m}\, Q(x)\, \eu^{\gamma x},
\qquad
m = \text{kegandaan } \gamma \text{ sebagai akar } \chi
\ (m = 0, 1 \text{ atau } 2),
\]
dengan $Q$ polinomial berderajat sama dengan $P$, yang koefisiennya
ditemukan lewat substitusi dan penyamaan. Untuk $f = K\cos\omega x$
(atau $\sin$), selesaikan dengan ruas kanan $K\eu^{\iu\omega x}$ lalu ambil
bagian realnya (masing-masing bagian imajinernya).
\end{method}

\begin{example}[Superposisi dalam praktik]\label{ex:b1:diffeq:superposition}
Selesaikan $y'' - y = \eu^{x} + 4$ pada $\R$. Homogennya: $\chi(r) = r^2
- 1$, dengan akar $\pm1$, jadi $y_h = \lambda\eu^x + \mu\eu^{-x}$. Belah
ruas kanannya lalu tangani setiap potongan dengan kotak metodenya.
\emph{Potongan $\eu^x$:} di sini $\gamma = 1$ adalah akar sederhana
$\chi$, jadi cobalah $y_1 = c\,x\,\eu^x$: maka $y_1'' - y_1 = c(x +
2)\eu^x - cx\eu^x = 2c\,\eu^x$, yang memberikan $c = \frac12$.
\emph{Potongan $4$:} di sini $\gamma = 0$ bukan akarnya; konstanta $y_2 =
-4$ memenuhi. Menurut superposisi
(\cref{thm:b1:diffeq:structure2}~(2)):
\[
y = \frac{x\,\eu^x}2 - 4 + \lambda\,\eu^x + \mu\,\eu^{-x},
\qquad (\lambda, \mu) \in \R^2 .
\]
Perhatikan bagaimana kedua potongan itu menuntut bentuk yang \emph{berbeda} ($m = 1$
berbanding $m = 0$): uji kegandaannya diterapkan pada masing-masing
eksponen secara terpisah, dan itulah seluruh maksud pembelahan
ruas kanannya sebelum menerka.
\end{example}

\begin{example}[Kaidah kegandaan dalam praktik]\label{ex:b1:diffeq:multiplicity}
Selesaikan $y'' + y' = x$ pada $\R$. Ruas kanannya adalah $P(x)\eu^{0
\cdot x}$ dengan $P(x) = x$, dan $\gamma = 0$ adalah akar \emph{sederhana}
$\chi(r) = r^2 + r = r(r + 1)$: jadi $m = 1$, dan terkaan
yang benar adalah $y_p = x\,(\alpha x + \beta) = \alpha x^2 + \beta x$, yang satu
derajat lebih tinggi daripada $P$. Dengan mensubstitusikannya:
\[
y_p'' + y_p' = 2\alpha + (2\alpha x + \beta)
= 2\alpha x + (2\alpha + \beta) ,
\]
dan penyamaan dengan $x$ memberikan $\alpha = \frac12$, $\beta = -1$:
jadi $y_p = \frac{x^2}2 - x$. Penyelesaian umumnya: $y = \frac{x^2}2 - x +
\lambda + \mu\,\eu^{-x}$. Seandainya kita menerka $y_p = \alpha x + \beta$
(dengan mengabaikan kegandaannya), substitusi akan memberikan $y_p'' + y_p'
= \alpha$, yaitu sebuah konstanta --- tak ada pilihan $\alpha, \beta$ yang dapat mencocoki
$x$, dan kegagalannya bersifat struktural: konstanta sudah memenuhi
persamaan homogennya, jadi ia tak terlihat oleh ruas kirinya.
Faktor $x^m$ ada justru untuk memanjat keluar dari ruang penyelesaian
homogennya.
\end{example}

\begin{remark}[Polis asuransi tiga puluh detik]\label{rem:b1:diffeq:check}
Setiap persamaan yang diselesaikan pada bab ini ditutup dengan pemeriksaan
substitusi, dan itu bukan hiasan. Perhitungan persamaan
diferensial merangkai banyak langkah kecil (sebuah antiturunan, sebuah aturan hasil kali,
dua konstanta), dan satu salah tanda merambat tanpa terlihat;
mensubstitusikan kembali rumus akhirnya ke dalam persamaan menangkap
hampir semuanya dengan ongkos satu kali penurunan.
Peliharalah refleks itu dalam tiga lapis: periksa \emph{penyelesaian
khusus} itu sendiri (bagian homogennya toh saling hapus), periksa
\emph{syarat awal}nya pada penyelesaian penuh, dan bila ada
parameter, periksa sebuah \emph{nilai merosot} (apakah
rumus untuk $\Omega$ yang umum menghasilkan jawaban yang sudah dikenal di
$\Omega = 0$?). Kebiasaan itu berongkos setengah menit; ia mengubah
``mungkin benar'' menjadi ``sudah terverifikasi''.
\end{remark}

\begin{remark}[Jebakan yang lazim]\label{rem:b1:diffeq:pitfalls}
\begin{enumerate}
  \item \emph{Normalkan lebih dulu.} Rumusnya mengandaikan persamaannya
        berbunyi $y' + a(x)y = b(x)$ --- dengan koefisien $1$ pada $y'$.
        Untuk $xy' - 2y = x^3$, bagilah dengan $x$ (pada selang
        yang menghindari $0$) sebelum mengenali $a$ dan $b$, seperti pada
        \cref{exo:b1:diffeq:2}.
  \item \emph{Satu konstanta per dimensi, yang ditetapkan pada akhirnya.} Penyelesaian
        umum orde satu memikul satu konstanta, dan yang orde
        dua memikul dua; syarat awalnya dikenakan pada
        penyelesaian yang \emph{lengkap} $y_p + y_h$, jangan pernah pada $y_h$
        saja --- mengenakannya sebelum menambahkan $y_p$ adalah kesalahan
        struktural yang paling sering terjadi.
  \item \emph{Perhatikan kegandaannya.} Terkaan penyelesaian khusus
        yang memenuhi persamaan homogen tak terlihat oleh
        ruas kirinya; faktor $x^m$ pada
        \cref{met:b1:diffeq:particular} bukanlah pilihan
        (\cref{ex:b1:diffeq:multiplicity}).
  \item \emph{Selangnya adalah bagian dari jawabannya.} Penyelesaian hidup
        pada selang yang di situ koefisiennya kontinu;
        merekatkannya melewati kesingularan dapat melahirkan konstanta palsu
        (\cref{exo:b1:diffeq:12}) atau merusak ketunggalannya. ``Selesaikan
        pada $\R^*$'' berarti dua masalah yang saling bebas.
\end{enumerate}
\end{remark}

\begin{example}[Pemaksaan di luar resonansi]\label{ex:b1:diffeq:offresonance}
Selesaikan $y'' + 4y = \sin x$ pada $\R$. Frekuensi alaminya $2$, frekuensi
pemaksanya $1$: karena $\iu$ \emph{bukan} akar $\chi(r) = r^2
+ 4$, kegandaannya $m = 0$ dan sinusoid biasa sudah cukup.
Dengan mencoba $y_p = \alpha\sin x$ (tanpa perlu kosinus: persamaannya tak
mempunyai suku $y'$, dan $\sin$ melahirkan kembali $\sin$):
\[
y_p'' + 4y_p = -\alpha\sin x + 4\alpha\sin x = 3\alpha\sin x ,
\]
jadi $\alpha = \frac13$ dan penyelesaian umumnya adalah
\[
y = \frac{\sin x}3 + \lambda\cos 2x + \mu\sin 2x .
\]
Setiap penyelesaiannya tetap terbatas: yakni superposisi dua osilasi
pada frekuensi $1$ (yang dipaksakan) dan $2$ (yang alami). Bandingkan dengan
contoh berikutnya, yang di situ pemaksaan \emph{pada} frekuensi alaminya
mengubah bentuk jawabannya sama sekali.
\end{example}

\begin{example}[Sebuah osilasi terpaksa]\label{ex:b1:diffeq:oscillation}
Selesaikan $y'' + y = \cos x$, $y(0) = 0$, $y'(0) = 0$.

\emph{Homogennya:} $\chi(r) = r^2 + 1$, dengan akar $\pm\iu$: jadi $y_h =
\lambda\cos x + \mu \sin x$.

\emph{Khususnya:} ruas kanannya $\Re(\eu^{\iu x})$ dengan $\gamma =
\iu$ akar sederhana $\chi$: cobalah $z_p = c\, x\, \eu^{\iu x}$
($c \in \C$). Maka $z_p'' + z_p = c\,(2\iu)\eu^{\iu x}$, yang sama dengan
$\eu^{\iu x}$ untuk $c = \frac{1}{2\iu} = -\frac\iu2$. Jadi $z_p =
-\frac{\iu}{2} x (\cos x + \iu \sin x)$ dan $y_p = \Re(z_p) =
\frac{x \sin x}{2}$.

\emph{Penyelesaian umumnya:} $y = \frac{x\sin x}{2} + \lambda\cos x +
\mu\sin x$. Syaratnya: $y(0) = \lambda = 0$; lalu $y' = \frac{\sin x +
x\cos x}{2} + \mu\cos x$, sehingga $y'(0) = \mu = 0$. Jawabannya: $y =
\frac{x\sin x}{2}$ --- yaitu osilasi yang amplitudonya tumbuh linear:
inilah gejala \emph{resonansi}\index{resonansi}, yang disebabkan oleh pemaksaan
sistem itu pada frekuensi alaminya.
\end{example}

\begin{omfigure}
\begin{tikzpicture}
\begin{axis}[omaxis, width=11cm, height=5.2cm,
    xmin=0, xmax=25, ymin=-13, ymax=13,
    xtick={0,5,10,15,20,25}, ytick=\empty]
  \addplot[omThm, thick, domain=0:25, samples=400] {x*sin(deg(x))/2};
  \addplot[dashed, gray, domain=0:25] {x/2};
  \addplot[dashed, gray, domain=0:25] {-x/2};
\end{axis}
\end{tikzpicture}

{\small \omterm{ex:b1:diffeq:oscillation}{Resonansi}: penyelesaian $y = \frac{x \sin x}{2}$ bagi $y'' + y =
\cos x$ berayun di antara garis $y = \pm\frac x2$ (garis putus-putus), dengan
amplitudo yang terus tumbuh.}
\end{omfigure}

\begin{remark}[Selingan: kelinearan adalah sebuah geometri]\label{rem:b1:diffeq:interlude}
Lihat kembali bentuk setiap \omterm{def:b1:logic:sets}{himpunan} penyelesaian pada bab ini: sebuah
penyelesaian istimewa ditambah sebuah ruang penyelesaian homogen dengan satu
konstanta bebas (orde satu) atau dua (orde dua). Bab tentang
aljabar linear (\cref{ch:b1:vspaces,ch:b1:findim,ch:b1:linmaps})
akan memasok kosakata yang tepat: \omterm{def:b1:logic:map}{pemetaan} $L(y) = y'' + ay' + by$
bersifat \emph{linear}, penyelesaian homogennya membentuk \emph{kernel}
$L$, yaitu ruang vektor yang \emph{dimensinya} sama dengan orde persamaannya
--- itulah isi jujur ungkapan ``satu konstanta per orde'' --- dan
\omterm{def:b1:logic:sets}{himpunan} penyelesaian $L(y) = f$ adalah sebuah \emph{subruang afin}, yaitu
geseran kernelnya. Bahkan \omterm{def:b1:logic:map}{pemetaan} Cauchy $(\lambda, \mu)
\mapsto (y(x_0), y'(x_0))$ pada \cref{thm:b1:diffeq:structure2} adalah
bijeksi linear antara dua bidang, yakni sistem $2
\times 2$ yang punya invers (\cref{ch:b1:matrices}). Tak ada satu pun pada bab ini
yang perlu dikerjakan ulang --- ia hanya perlu dinamai ulang, dan
penamaan ulang itu adalah pemanasan terbaik bagi aljabar linear:
setiap definisi abstrak di sana sudah mencari nafkahnya
di sini.
\end{remark}

\begin{remark}[Di mana bab ini dipakai]\label{rem:b1:diffeq:whereused}
Teorema strukturnya --- bahwa penyelesaian $(E)$ berupa ``satu penyelesaian
khusus ditambah penyelesaian $(H)$'' --- adalah kemunculan pertama
sebuah pola yang akan dinamai \cref{ch:b1:vspaces,ch:b1:linmaps}:
\omterm{def:b1:logic:sets}{himpunan} penyelesaian $(H)$ adalah \emph{kernel} \omterm{def:b1:logic:map}{pemetaan} linear $y
\mapsto y'' + ay' + by$, dan \omterm{def:b1:logic:sets}{himpunan} penyelesaian $(E)$ adalah
geseran afinnya. \omterm{def:b1:diffeq:linear2}{Polinomial karakteristiknya} muncul kembali sebagai
\omterm{def:b1:diffeq:linear2}{polinomial karakteristik} sebuah matriks pada \cref{ch:b1:matrices}:
persamaan orde dua sesungguhnya sistem orde satu $2 \times 2$ yang
menyamar, sudut pandang yang disistematiskan jilid Tahun ke-2. Integral
yang dituntut \omterm{thm:b1:diffeq:voc}{variasi konstanta} dipasok oleh
\cref{ch:b1:integration}, dan soal akhir pekan di bawah --- yaitu
osilator teredam yang dipaksa --- adalah kasus model bagi setiap
pertanyaan osilasi dalam sains, dari rangkaian listrik sampai jembatan
gantung.
\end{remark}

\section{Latihan}

\begin{exercise}[$\star$]\label{exo:b1:diffeq:1}
Selesaikan pada $\R$: $\;y' + 2y = \eu^{3x}$; lalu \omterm{thm:b1:diffeq:voc}{masalah Cauchy} $y(0) =
1$.
\end{exercise}

\begin{exercise}[$\star$]\label{exo:b1:diffeq:2}
Selesaikan pada $\intoo{0}{+\infty}$: $\;x y' - 2y = x^3$ \emph{(bawalah
persamaannya ke bentuk yang ternormalkan lebih dulu)}.
\end{exercise}

\begin{exercise}[$\star$]\label{exo:b1:diffeq:3}
Selesaikan pada $\R$, dengan memberikan penyelesaian umum yang real:
$\;y'' - 3y' + 2y = 0$; $\;y'' + 4y' + 4y = 0$; $\;y'' - 2y' + 5y =
0$.
\end{exercise}

\begin{exercise}[$\star$]\label{exo:b1:diffeq:4}
Selesaikan $y'' - y = x^2$ pada $\R$, lalu \omterm{thm:b1:diffeq:voc}{masalah Cauchy} $y(0) = 0$,
$y'(0) = 1$.
\end{exercise}

\begin{exercise}[$\star\star$]\label{exo:b1:diffeq:5}
Selesaikan pada $\intoo{-\frac\pi2}{\frac\pi2}$:
$\;y' + y\tan x = \sin 2x$.
\end{exercise}

\begin{exercise}[$\star\star$]\label{exo:b1:diffeq:6}
Selesaikan $y'' - 4y' + 3y = (2x + 1)\,\eu^{x}$ pada $\R$. \emph{(Perhatikan
kegandaannya: apakah $1$ akar \omterm{def:b1:diffeq:linear2}{polinomial karakteristiknya}?)}
\end{exercise}

\begin{exercise}[$\star\star$]\label{exo:b1:diffeq:7}
Selesaikan $y'' + 4y = \sin 2x + x$ pada $\R$ \emph{(superposisi; tangani
setiap ruas kanannya secara terpisah)}.
\end{exercise}

\begin{exercise}[$\star\star$]\label{exo:b1:diffeq:8}
Secangkir kopi bersuhu $T_0 = 80\,^\circ$C berada di ruangan bersuhu
$20\,^\circ$C. Hukum pendinginan Newton menyatakan $T' = -k\,(T - 20)$ dengan
$k > 0$. Selesaikan untuk $T(t)$, lalu bila diketahui kopinya bersuhu
$50\,^\circ$C setelah $10$ menit, carilah kapan ia mencapai
$25\,^\circ$C.
\end{exercise}

\begin{exercise}[$\star\star\star$]\label{exo:b1:diffeq:9}
(Osilator teredam) Untuk $\varepsilon \geq 0$, tinjau
$y'' + 2\varepsilon y' + y = 0$.
\begin{enumerate}
  \item Selesaikan untuk $\varepsilon \in \intco{0}{1}$, $\varepsilon = 1$,
        dan $\varepsilon > 1$.
  \item Tunjukkan bahwa untuk $\varepsilon > 0$ setiap penyelesaiannya menuju $0$
        di $+\infty$, dan bahwa untuk $\varepsilon = 0$ penyelesaian yang tak
        nol tidak demikian.
  \item Untuk $\varepsilon \in \intoo{0}{1}$, tunjukkan bahwa nol sebuah
        penyelesaian tak nol berjarak teratur, dengan celah
        $\frac{\pi}{\sqrt{1 - \varepsilon^2}}$.
\end{enumerate}
\end{exercise}

\begin{exercise}[$\star\star\star$]\label{exo:b1:diffeq:10}
Carilah semua fungsi $f \colon \R \to \R$ yang dapat diturunkan dua kali dan
memenuhi
\[
\forall x, y \in \R, \qquad f(x + y) + f(x - y) = 2 f(x) f(y),
\]
dengan $f(0) \ne 0$ dan $f$ tidak konstan.
\emph{Petunjuk: tetapkan $y$, lalu turunkan dua kali terhadap $x$ di $0$;
tunjukkan $f(0) = 1$ dan $f'' = c f$ untuk suatu konstanta $c$; lalu selesaikan
menurut tanda $c$ dan periksa penyelesaian mana yang memenuhi
persamaan fungsionalnya.}
\end{exercise}

\begin{exercise}[$\star\star$]\label{exo:b1:diffeq:11}
(Persamaan Euler) Selesaikan $x^2 y'' - x y' + y = 0$ pada
$\intoo{0}{+\infty}$. \emph{Petunjuk: tulis $z(t) = y(\eu^t)$, yakni
substitusikan $x = \eu^t$, lalu tunjukkan bahwa $z$ memenuhi persamaan
linear berkoefisien konstan.}
\end{exercise}

\begin{exercise}[$\star\star\star$]\label{exo:b1:diffeq:12}
Tinjau persamaan $x\,y' = 2y$ pada seluruh garis real, dengan
fungsi $y \colon \R \to \R$ yang dapat diturunkan sebagai variabelnya.
\begin{enumerate}
  \item Selesaikan pada $\intoo{0}{+\infty}$ dan pada $\intoo{-\infty}{0}$.
  \item Tunjukkan bahwa untuk konstanta $a, b \in \R$ yang \emph{mana pun},
        fungsi yang sama dengan $ax^2$ untuk $x \geq 0$ dan dengan $bx^2$ untuk
        $x < 0$ dapat diturunkan pada $\R$ dan memenuhi persamaannya
        di mana-mana.
  \item Simpulkan bahwa \omterm{def:b1:logic:sets}{himpunan} penyelesaiannya pada $\R$ adalah keluarga
        berparameter dua, lalu jelaskan mengapa ini tidak bertentangan dengan
        ketunggalan pada \cref{thm:b1:diffeq:voc}.
\end{enumerate}
\end{exercise}

\section{Soal: Osilator teredam yang dipaksa}

\begin{problem}\label{pb:b1:diffeq:1}
Satu persamaan mengatur sebuah massa pada pegas di dalam medium kental,
muatan pada rangkaian RLC, dan gedung yang bergoyang tertiup angin:
\[
(E_\Omega)\colon\quad
x'' + 2\lambda x' + \omega_0^2\,x = A\cos(\Omega t),
\]
dengan $\lambda \geq 0$ redamannya, $\omega_0 > 0$ frekuensi
alaminya, serta $A > 0$, $\Omega > 0$ amplitudo dan frekuensi
pemaksanya. Soal ini menyarikan perilakunya selengkapnya: peluruhan
peralihannya, keadaan tunak periodik yang tunggal, kurva \omterm{ex:b1:diffeq:oscillation}{resonansi}
beserta ketajamannya (yaitu \emph{\omterm{pb:b1:diffeq:1}{faktor kualitas}}), \omterm{pb:b1:diffeq:1}{layangan} pada
kasus tanpa redaman, dan neraca energi yang menopang
osilasinya.\index{resonansi}\index{osilator}
Kecuali dinyatakan lain, $0 < \lambda < \omega_0$ (yaitu regim redaman
lemah) dan kita tulis $\omega_d = \sqrt{\omega_0^2 - \lambda^2}$.

\medskip
\noindent\textbf{Bagian I --- Osilator bebas.} Di sini $A = 0$.
\begin{enumerate}
  \item Selesaikan persamaan homogennya $(H)$ untuk $0 < \lambda <
        \omega_0$, dan untuk $\lambda = 0$. (Regim $\lambda
        \geq \omega_0$ sudah ditangani pada \cref{exo:b1:diffeq:9};
        kutiplah dari sana.)
  \item Tunjukkan bahwa untuk setiap $\lambda > 0$, semua penyelesaian $(H)$
        menuju $0$ di $+\infty$ --- pada ketiga regimnya.
  \item Definisikan energinya $\mathcal E(t) = \frac12 x'(t)^2 +
        \frac12\omega_0^2\,x(t)^2$ sepanjang sebuah penyelesaian $(H)$. Tunjukkan
        $\mathcal E'(t) = -2\lambda\,x'(t)^2 \leq 0$, lalu simpulkan
        (tanpa menyelesaikan apa pun) bahwa \omterm{thm:b1:diffeq:voc}{masalah Cauchy}
        ``$(H)$, $x(t_0) = x'(t_0) = 0$'' hanya mempunyai penyelesaian
        nol, untuk setiap $\lambda \geq 0$.
  \item Untuk $0 < \lambda < \omega_0$, tulis penyelesaian yang tak nol sebagai
        $x(t) = R\,\eu^{-\lambda t}\cos(\omega_d t - \varphi)$ lalu
        misalkan $T_d = \frac{2\pi}{\omega_d}$ pseudoperiodenya.
        Tunjukkan bahwa $x(t + T_d) = \eu^{-\lambda T_d}\,x(t)$: jadi setiap
        ayunan adalah ayunan sebelumnya yang menyusut dengan faktor tetap
        $\eu^{-\delta}$, $\delta = \frac{2\pi\lambda}{\omega_d}$
        (yaitu \emph{dekremen logaritmiknya}). Hitung $\delta$ untuk
        $\omega_0 = 1$, $\lambda = 0.1$.
  \item Definisikan \emph{\omterm{pb:b1:diffeq:1}{faktor kualitas}} $Q =
        \dfrac{\omega_0}{2\lambda}$\index{faktor kualitas}. Tunjukkan
        bahwa setelah waktu $\frac1\lambda$ (yaitu satu kali pelipatan
        amplitudo sebesar $\eu$), osilatornya sudah menyelesaikan
        $\frac{\omega_d}{2\pi\lambda}$ pseudoperiode, yang untuk
        redaman lemah ($\lambda \ll \omega_0$) kira-kira sama dengan
        $\frac Q\pi$: jadi \omterm{pb:b1:diffeq:1}{faktor kualitas} mencacah, sampai faktor $\pi$,
        banyaknya osilasi yang bertahan sebelum amplitudonya meluruh sebesar
        $\eu$.
\end{enumerate}

\medskip
\noindent\textbf{Bagian II --- Keadaan tunak.} Sekarang $A > 0$ dan
$\lambda > 0$.
\begin{enumerate}[resume]
  \item Carilah penyelesaian khusus sebagai bagian real
        $z\,\eu^{\iu\Omega t}$ dengan $z \in \C$
        (\cref{met:b1:diffeq:particular}). Tunjukkan bahwa ini berhasil
        dengan
        \[
        z = \frac{A}{\omega_0^2 - \Omega^2 + 2\iu\lambda\Omega} .
        \]
  \item Simpulkan keadaan tunaknya dalam bentuk amplitudo--fase:
        $x_p(t) = R(\Omega)\cos\bigl(\Omega t -
        \varphi(\Omega)\bigr)$ dengan
        \[
        R(\Omega) = \frac{A}{\sqrt{(\omega_0^2 - \Omega^2)^2 +
        4\lambda^2\Omega^2}},
        \qquad
        \tan\varphi = \frac{2\lambda\Omega}{\omega_0^2 -
        \Omega^2},
        \quad \varphi \in \intoo0\pi .
        \]
  \item Tafsirkan kedua regim ekstremnya: hitung limit
        $R$ dan $\varphi$ ketika $\Omega \to 0^+$ (tanggapan
        kuasi-statis $A/\omega_0^2$, dengan fase $0$) dan ketika $\Omega \to
        +\infty$ ($R \sim A/\Omega^2 \to 0$, dengan fase $\to \pi$: jadi
        massanya bergerak berlawanan dengan pemaksaan yang terlalu cepat).
  \item Tunjukkan bahwa \emph{setiap} penyelesaian $(E_\Omega)$ berupa $x_p$
        ditambah sebuah penyelesaian $(H)$, sehingga ia konvergen ke keadaan
        tunak $x_p$ ketika $t \to +\infty$, apa pun syarat
        awalnya: setelah peralihannya mati, osilatornya tak
        mempunyai ingatan tentang bagaimana ia dimulai.
  \item Tunjukkan bahwa $x_p$ adalah \emph{satu-satunya} penyelesaian periodik
        $(E_\Omega)$.
  \item Kerjakan satu \omterm{thm:b1:diffeq:voc}{masalah Cauchy} sampai tuntas: untuk $x'' + 2x' + 2x =
        \cos t$ dengan $x(0) = x'(0) = 0$, tunjukkan bahwa penyelesaiannya
        \[
        x(t) = \frac{\cos t + 2\sin t}5
        - \eu^{-t}\,\frac{\cos t + 3\sin t}5 ,
        \]
        lalu kenali bagian peralihan dan bagian tunaknya.
\end{enumerate}

\medskip
\noindent\textbf{Bagian III --- Kurva \omterm{ex:b1:diffeq:oscillation}{resonansi}.} Telaah
$\Omega \mapsto R(\Omega)$ pada $\intoo0{+\infty}$.
\begin{enumerate}[resume]
  \item Dengan menulis $u = \Omega^2$ dan $g(u) = (\omega_0^2 - u)^2 +
        4\lambda^2 u$, tunjukkan: jika $2\lambda^2 < \omega_0^2$, maka $R$
        mencapai maksimum tegas pada \emph{frekuensi \omterm{ex:b1:diffeq:oscillation}{resonansi}}
        $\Omega_r = \sqrt{\omega_0^2 - 2\lambda^2}$, dengan
        \[
        R_{\max} = R(\Omega_r)
        = \frac{A}{2\lambda\sqrt{\omega_0^2 - \lambda^2}} .
        \]
  \item Tunjukkan bahwa $\dfrac{R_{\max}}{R(0)} = Q\,\bigl(1 -
        \tfrac{\lambda^2}{\omega_0^2}\bigr)^{-1/2} \geq Q$: jadi pada
        \omterm{ex:b1:diffeq:oscillation}{resonansi}, pemaksaannya diperkuat sebesar (pada dasarnya) faktor
        kualitasnya.
  \item Buktikan bahwa amplitudo \emph{kecepatan} $V(\Omega) =
        \Omega\,R(\Omega)$ maksimum tepat di $\Omega =
        \omega_0$ (bukan di $\Omega_r$), dan bahwa fase di sana adalah
        $\varphi(\omega_0) = \frac\pi2$: jadi di $\Omega = \omega_0$
        kecepatannya tepat sefase dengan gayanya.
  \item (Lebar pita) Selesaikan $g(u) = 2\,g(u_r)$ secara eksak, dengan $u_r =
        \omega_0^2 - 2\lambda^2$, lalu simpulkan bahwa kedua
        frekuensi $\Omega_\pm$ yang di sana $R = R_{\max}/\sqrt2$
        memenuhi $\Omega_+^2 - \Omega_-^2 =
        4\lambda\sqrt{\omega_0^2 - \lambda^2}$; lalu simpulkan bahwa untuk
        redaman lemah lebar pitanya $\Omega_+ - \Omega_-
        \approx 2\lambda$, yakni $Q \approx
        \frac{\omega_0}{\Omega_+ - \Omega_-}$: jadi puncak \omterm{ex:b1:diffeq:oscillation}{resonansi}
        yang tajam adalah sistem ber-$Q$ tinggi.
  \item Potret numerik untuk $\omega_0 = 1$, $\lambda = 0.05$
        ($Q = 10$), $A = 1$: hitung $\Omega_r$, $R_{\max}$,
        tanggapan statis $R(0)$, dan lebar pitanya secara hampiran.
  \item Tunjukkan bahwa jika $2\lambda^2 \geq \omega_0^2$, maka $R$
        turun tegas pada $\intoo0{+\infty}$: jadi sistem yang teredam
        berat sama sekali tak mempunyai puncak \omterm{ex:b1:diffeq:oscillation}{resonansi}.
\end{enumerate}

\medskip
\noindent\textbf{Bagian IV --- Tanpa redaman: \omterm{pb:b1:diffeq:1}{layangan} dan \omterm{ex:b1:diffeq:oscillation}{resonansi}.}
Di sini $\lambda = 0$.
\begin{enumerate}[resume]
  \item Untuk $\Omega \neq \omega_0$, carilah penyelesaian umum
        $x'' + \omega_0^2 x = A\cos(\Omega t)$.
  \item Selesaikan \omterm{thm:b1:diffeq:voc}{masalah Cauchy} $x(0) = x'(0) = 0$ lalu ubahlah
        jawabannya menjadi bentuk hasil kali
        \[
        x(t) = \frac{2A}{\omega_0^2 - \Omega^2}\,
        \sin\Bigl(\frac{(\omega_0 - \Omega)t}2\Bigr)
        \sin\Bigl(\frac{(\omega_0 + \Omega)t}2\Bigr) .
        \]
  \item Untuk $\Omega$ yang dekat dengan $\omega_0$, bacalah hasil kali itu sebagai
        osilasi cepat pada frekuensi $\frac{\omega_0 + \Omega}2$
        yang termodulasi oleh selubung lambat pada frekuensi
        $\frac{\abs{\omega_0 - \Omega}}2$: inilah \emph{layangan}
        \index{layangan}. Berikan periode selubungnya dan
        amplitudo maksimumnya, lalu catat bagaimana keduanya meledak ketika $\Omega
        \to \omega_0$.
  \item Tetapkan $t$ lalu ambil $\Omega \to \omega_0$ pada rumus
        pertanyaan 19: tunjukkan bahwa limitnya adalah
        \[
        x_\infty(t) = \frac{A\,t\,\sin(\omega_0 t)}{2\omega_0} ,
        \]
        lalu periksa langsung bahwa $x_\infty$ memenuhi persamaan
        resonannya $x'' + \omega_0^2 x = A\cos(\omega_0 t)$ dengan
        $x(0) = x'(0) = 0$ (bandingkan
        \cref{ex:b1:diffeq:oscillation}): jadi \omterm{ex:b1:diffeq:oscillation}{resonansi} adalah limit
        \omterm{pb:b1:diffeq:1}{layangan} yang kian lambat dan kian besar.
  \item Pertentangkan kedua nasib \omterm{ex:b1:diffeq:oscillation}{resonansi}: pertumbuhan linear
        $\frac{At}{2\omega_0}$ tanpa redaman, berbanding kejenuhan
        pada $R_{\max} \approx Q\,\frac{A}{\omega_0^2}$ dengan redaman
        lemah. Dalam satu kalimat: mekanisme fisis apa yang mengubah
        yang pertama menjadi yang kedua?
\end{enumerate}

\medskip
\noindent\textbf{Bagian V --- Neraca energi dan sintesis.}
\begin{enumerate}[resume]
  \item Pada keadaan tunak Bagian II, hitunglah rata-rata sepanjang
        satu periode $\frac{2\pi}\Omega$ dari (a) daya yang disuntikkan
        oleh pemaksaannya, $P_{\mathrm{in}}(t) = A\cos(\Omega t)
        \cdot x_p'(t)$, dan (b) daya yang dilesapkan oleh
        redamannya, $P_{\mathrm{diss}}(t) = 2\lambda\,x_p'(t)^2$.
        Tunjukkan bahwa kedua rata-ratanya sama dengan $\lambda\,R^2\Omega^2$: jadi
        pemaksaannya memasok persis apa yang dibakar redamannya --- dan
        itulah sebabnya keadaan tunak itu tunak.
  \item Di mana persisnya soal ini memakai: (i) teorema struktur
        \cref{thm:b1:diffeq:structure2}; (ii) metode eksponensial
        kompleks; (iii) telaah fungsi variabel real dalam gaya
        \cref{ch:b1:functions}? Satu kalimat untuk masing-masing.
  \item Sintesis: gambarkan peta perilaku lengkap $(E_\Omega)$
        --- yang bebas berbanding yang terpaksa, yang teredam berbanding tanpa redaman, peran
        $Q$ sebagai satu-satunya tombol tanpa dimensi yang menyetel tinggi puncak,
        lebar pita dan umur peralihannya --- lalu sebutkan
        tempat kisahnya berlanjut: sistem orde satu $2 \times 2$
        (\cref{ch:b1:matrices} dan jilid Tahun ke-2) dan
        penguraian pemaksaan periodik yang umum menjadi
        sinusoid (deret Fourier, pada jilid Tahun ke-3), yang
        untuknya kasus sinusoidal soal ini menjadi batu bangunan
        yang mendasar.

\end{enumerate}
\end{problem}
